All praise and gratitude be to Allah SWT for His abundant grace, blessings, and guidance, which have enabled the author to successfully complete this final project in the form of a thesis titled "Development of RayCastED: A Lightweight Contour-based Nuclei Instance Segmentation for H\&E Histopathology Slides"

During the preparation and writing of this thesis, the author has received invaluable guidance, assistance, and support from various parties. Therefore, on this occasion, the author would like to express the deepest gratitude to:
\begin{enumerate}
	\item Ir. Hanung Adi Nugroho, S.T., M.E., Ph.D., IPM., as the Head of the Department of Electrical Engineering and Information Technology.
	
	\item Syukron Abu Ishaq Alfarozi, S.T., Ph.D., Head of the Undergraduate Program in Biomedical Engineering.
		
	\item Ir. Ridwan Wicaksono, S.T., M.Eng., Ph.D., as the first supervisor, for his guidance, support, and valuable advice throughout the research and thesis writing process.

	\item dr. Paranita Ferronika, Ph.D., Sp.PA., Subsp,KA(K), as the second supervisor, for her guidance, support, and valuable advice throughout the research and thesis writing process.

	\item My beloved parents, Palupi Dewi and Irwan Taufiq Ritonga, and my brother, Roshan Rakin Ritonga, for their unwavering support, prayers, encouragement, and faith throughout my academic journey.

	\item All dearest friends, Kayana Anindya Azaria, Gavind Muhammad Pramahita, Raihan Fauzy Ferdyansyah, Caroline Olivia, Azfazaki Hakimi, Faisol Rukmana A., Ganisya Rihma A., Melody Nadaa F., Ariq Fadhilah J., and all friends in the Biomedical Engineering Undergraduate Program, for their support, encouragement, and companionship throughout this journey.

	\item Fellow lab colleagues, Yahya Ayas F., Ali Mochtar A. S., Chenaniah, Salsabila Rannandita P., and Polar Osaka for their support, collaboration, and valuable discussions during the research process.

	\item The friend group(s), P Info and Pejuang RL, for all the companionship, support, and memorable moments shared during this journey.

\end{enumerate}

Finally, the author hopes that this thesis will be beneficial and provide meaningful contributions to all of us, Aamiin.

%\begin{flushright}
%	\begin{tabular}{c}
%		Yogyakarta, 13 Maret 2017 \\
%		\vspace{1cm} \\
%		Canggih
%	\end{tabular}
%\end{flushright}
